% ---------------------------------------------------------------------------
% Shared TikZ styles for paper figures.
%
% Drop this in the preamble, then draw every diagram with these styles. One edit
% here restyles every figure in the paper, which is the whole point: consistency
% across figures is what makes a paper look professionally produced.
%
% Colors are the Okabe-Ito colorblind-safe palette. Meaning is never carried by
% color alone: use `pfig/box` for context and `pfig/new` for what the paper
% contributes, and the contrast survives a grayscale print.
%
% Requires: \usepackage{tikz} and the libraries loaded below. Builds with stock
% TeX Live, no external binaries.
% ---------------------------------------------------------------------------

\usetikzlibrary{arrows.meta,positioning,fit,backgrounds,calc,shapes.geometric}

% Okabe-Ito palette ---------------------------------------------------------
\definecolor{okBlue}{HTML}{0072B2}
\definecolor{okSky}{HTML}{56B4E9}
\definecolor{okGreen}{HTML}{009E73}
\definecolor{okOrange}{HTML}{E69F00}
\definecolor{okVermillion}{HTML}{D55E00}
\definecolor{okPurple}{HTML}{CC79A7}
\definecolor{okYellow}{HTML}{F0E442}
\definecolor{pfigGray}{HTML}{6E6E6E}
\definecolor{pfigLight}{HTML}{EDEDED}

\tikzset{
  % Base node: context components. Deliberately quiet.
  pfig/box/.style={
    draw=pfigGray, fill=white, line width=0.6pt,
    rounded corners=2pt, inner sep=4pt, align=center,
    font=\footnotesize, minimum height=7mm, text=black
  },
  % The contribution. One accent, used sparingly: if everything is emphasized,
  % nothing is.
  pfig/new/.style={
    pfig/box, draw=okBlue, fill=okSky!18, line width=1pt, font=\footnotesize\bfseries
  },
  % Secondary emphasis, when a second component genuinely needs to stand out.
  pfig/alt/.style={
    pfig/box, draw=okOrange, fill=okOrange!12, line width=0.9pt
  },
  % External input or user.
  pfig/io/.style={
    pfig/box, rounded corners=6pt, fill=pfigLight, draw=pfigGray
  },
  % Data store.
  pfig/store/.style={
    pfig/box, cylinder, shape border rotate=90, aspect=0.25,
    inner sep=3pt, minimum width=16mm
  },
  % Decision or gate.
  pfig/gate/.style={
    pfig/box, diamond, aspect=1.8, inner sep=1pt, minimum height=6mm
  },
  % Swimlane or grouping background.
  pfig/lane/.style={
    draw=pfigGray!50, dashed, line width=0.5pt, rounded corners=3pt,
    fill=pfigLight!45, inner sep=5pt
  },
  % Arrows. Solid for the main path, dashed for feedback or optional flow.
  pfig/arrow/.style={
    -{Stealth[length=2.4mm,width=1.6mm]}, line width=0.7pt, draw=black
  },
  pfig/arrowD/.style={
    pfig/arrow, dashed, draw=pfigGray
  },
  pfig/arrowNew/.style={
    pfig/arrow, draw=okBlue, line width=0.9pt
  },
  % Edge labels: small, never smaller than 7pt at final size.
  pfig/label/.style={
    font=\scriptsize, text=black, inner sep=1.5pt, fill=white, fill opacity=0.85,
    text opacity=1
  },
  % Section caption inside a figure.
  pfig/title/.style={
    font=\footnotesize\bfseries, text=black
  },
}

% ---------------------------------------------------------------------------
% Example. Delete once you have your own figures.
%
% \begin{figure}[t]
%   \centering
%   \begin{tikzpicture}[node distance=6mm and 9mm]
%     \node[pfig/io]  (in)   {Workload};
%     \node[pfig/box, right=of in]  (enc) {Encoder};
%     \node[pfig/new, right=of enc] (ours){Verifier-guided\\memory};
%     \node[pfig/box, right=of ours](dec) {Decoder};
%     \node[pfig/io,  right=of dec] (out) {Answer};
%     \draw[pfig/arrow]    (in)  -- (enc);
%     \draw[pfig/arrowNew] (enc) -- node[pfig/label]{states} (ours);
%     \draw[pfig/arrowNew] (ours)-- (dec);
%     \draw[pfig/arrow]    (dec) -- (out);
%     \draw[pfig/arrowD]   (dec.south) to[out=-90,in=-90] node[pfig/label]{retry} (ours.south);
%   \end{tikzpicture}
%   \caption{...}
% \end{figure}
% ---------------------------------------------------------------------------
